\documentclass[11pt,a4paper]{ctexart}

\usepackage[a4paper,margin=17mm,headheight=23pt,footskip=10mm]{geometry}
\usepackage{amsmath,amssymb,bm,mathtools}
\usepackage{array,booktabs,tabularx,colortbl}
\usepackage{tikz}
\usetikzlibrary{arrows.meta,positioning,calc,fit}
\usepackage{fancyhdr}
\usepackage{xcolor}

\setCJKmainfont[AutoFakeBold=2]{SimSun}
\setCJKsansfont[AutoFakeBold=2]{SimSun}
\setmainfont{Times New Roman}
\setsansfont{Times New Roman}
\setmonofont{Consolas}

\definecolor{paperblack}{HTML}{FFFFFF}
\definecolor{panel}{HTML}{F7F9FC}
\definecolor{paneltwo}{HTML}{F1F5F9}
\definecolor{whiteish}{HTML}{111827}
\definecolor{muted}{HTML}{5B6573}
\definecolor{cyan}{HTML}{007F7B}
\definecolor{blue}{HTML}{1E5A91}
\definecolor{orange}{HTML}{C65D17}
\definecolor{green}{HTML}{2F855A}
\definecolor{grid}{HTML}{CBD5E1}
\definecolor{redsoft}{HTML}{B94343}

\pagecolor{paperblack}
\color{whiteish}
\arrayrulecolor{grid}
\setlength{\parindent}{0pt}
\setlength{\parskip}{5pt}
\renewcommand{\baselinestretch}{1.15}
\emergencystretch=2em

\pagestyle{fancy}
\fancyhf{}
\fancyfoot[C]{\color{muted}\small\thepage}
\renewcommand{\headrulewidth}{0.4pt}
\renewcommand{\headrule}{\hbox to\headwidth{\color{whiteish}\leaders\hrule height \headrulewidth\hfill}}

\newcommand{\vct}[1]{\bm{#1}}
\newcommand{\enc}{\operatorname{Enc}}
\newcommand{\reshape}{\operatorname{reshape}}
\newcommand{\shape}[1]{\textcolor{orange}{\ensuremath{#1}}}
\newcommand{\term}[1]{\textcolor{cyan}{\ensuremath{#1}}}
\newcommand{\query}[1]{\textcolor{blue}{\ensuremath{#1}}}
\newcommand{\cand}[1]{\textcolor{cyan}{\ensuremath{#1}}}

\newcommand{\docsection}[2]{%
  \vspace*{1mm}
  {\color{whiteish}\bfseries\Large #2}\par
  \vspace{1mm}{\color{whiteish}\rule{\linewidth}{0.8pt}}\vspace{2mm}
}

\newcommand{\keybox}[3][0.93\linewidth]{%
  \begin{center}
  \setlength{\fboxsep}{7pt}%
  \fcolorbox{whiteish}{panel}{%
    \begin{minipage}{#1}
    #3
    \end{minipage}}
  \end{center}
}

\newcommand{\smallbox}[2]{%
  \setlength{\fboxsep}{5pt}%
  \fcolorbox{whiteish}{paneltwo}{#2}%
}

\newcommand{\slot}[3]{%
  \tikz[baseline=-0.65ex]\node[draw=paperblack,fill=#1,text=white,
  minimum width=#2,minimum height=5.2mm,inner sep=1pt,font=\bfseries\scriptsize]{#3};%
}

\begin{document}

% ========================================================================
\docsection{cyan}{0. 问题定义与计算目标}

设 A 一次提交 $m$ 个查询，每个查询有 $d$ 维特征；B 有 $L$ 个候选。
把查询和候选分别按行排列：
\[
Q=\begin{bmatrix}q_0^{\mathsf T}\\[-1mm]\vdots\\q_{m-1}^{\mathsf T}\end{bmatrix}
\in\mathbb R^{m\times d},
\qquad
C=\begin{bmatrix}c_0^{\mathsf T}\\[-1mm]\vdots\\c_{L-1}^{\mathsf T}\end{bmatrix}
\in\mathbb R^{L\times d}.
\]

理想输出不是一个数，而是一张 $m\times L$ 的分数表：
\[
S=QC^{\mathsf T}\in\mathbb R^{m\times L},
\qquad
S_{q,t}=\langle q_q,c_t\rangle
=\sum_{f=0}^{d-1}Q_{q,f}C_{t,f}.
\]

\keybox{cyan}{
\textbf{打包问题。}
CKKS 的一个密文只有一维槽位。我们要构造一个一维索引，
使每个槽位唯一代表一个二维组合 $(q,t)$，并保证这个槽位内的逐特征乘加恰好得到 $S_{q,t}$。
}

为避免“维度在推导中悄悄变化”，先固定所有符号：

\begin{center}
\renewcommand{\arraystretch}{1.35}
\begin{tabularx}{0.96\linewidth}{>{\color{cyan}\bfseries}l >{\centering\arraybackslash}p{28mm} X}
\toprule
对象 & 形状 & 含义 \\
\midrule
$Q$ & $m\times d$ & A 的查询特征矩阵 \\
$C$ & $L\times d$ & B 的完整候选库 \\
$S_{\rm slots}$ & 标量 & 一个 CKKS 密文的可用槽位数，当前为 $4096$ \\
$T=\lfloor S_{\rm slots}/m\rfloor$ & 标量 & 每个查询在同一密文中可分到的候选位置数 \\
$C^{(j)}$ & $T\times d$ & 第 $j$ 个候选 tile；最后一块不足 $T$ 时补零 \\
$A^{\rm pack},B^{(j)}_{\rm pack}$ & $mT\times d$ & 查询与候选在加密前完成对齐后的二维矩阵 \\
\bottomrule
\end{tabularx}
\end{center}

\textcolor{muted}{后面始终使用 query-major 顺序：先放完 $q_0$ 的 $T$ 个候选位置，再放 $q_1$ 的 $T$ 个候选位置。}

\newpage

% ========================================================================
\docsection{blue}{1. 修改前的逐候选轮询计算}

旧方案只沿“查询方向”使用槽位。对每个特征 $f$，A 加密长度为 $m$ 的列向量：
\[
\vct x_f=Q_{:,f}\in\mathbb R^m,
\qquad \enc(\vct x_f).
\]

当 B 处理第 $t$ 个候选 $c_t$ 时，它把同一个候选特征值广播到全部查询槽位：
\[
\vct y_f^{(t)}=C_{t,f}\vct 1_m
=[C_{t,f},\ldots,C_{t,f}]^{\mathsf T}\in\mathbb R^m.
\]

于是一次逐槽乘加只能生成候选 $c_t$ 对应的一列：
\[
\vct z^{(t)}
=\sum_{f=0}^{d-1}\enc(\vct x_f)\odot\vct y_f^{(t)}
=\enc\!\left(
\begin{bmatrix}
\langle q_0,c_t\rangle\\[-1mm]
\vdots\\[-1mm]
\langle q_{m-1},c_t\rangle
\end{bmatrix}\right)
\in\enc(\mathbb R^m).
\]

\keybox{blue}{
\textbf{形状结论。}
$\vct x_f$ 和 $\vct y_f^{(t)}$ 都只有 $m$ 个有效位置，所以结果也只有 $m$ 个分数；
这 $m$ 个分数共享同一个候选编号 $t$。要得到完整 $m\times L$ 矩阵，只能令 $t=0,1,\ldots,L-1$ 逐列轮询。
}

用 $m=2,d=2$ 的例子看得最清楚：
\[
q_0=[1,2],\quad q_1=[3,4],\qquad c_0=[10,20].
\]

\begin{center}
\begin{tikzpicture}[node distance=8mm,>=Stealth]
\node (q) [draw=blue,rounded corners,fill=panel,inner sep=7pt,align=center]
{$\enc([1,3])$\hspace{3mm}与\hspace{3mm}$\enc([2,4])$\\[-1mm]
\textcolor{muted}{查询特征列，长度 $m=2$}};
\node (c) [right=of q,draw=cyan,rounded corners,fill=panel,inner sep=7pt,align=center]
{$[10,10]$\hspace{3mm}与\hspace{3mm}$[20,20]$\\[-1mm]
\textcolor{muted}{只广播候选 $c_0$}};
\node (r) [right=of c,draw=orange,rounded corners,fill=panel,inner sep=7pt,align=center]
{$\enc([50,110])$\\[-1mm]
\textcolor{muted}{只得到第 $0$ 列}};
\draw[->,line width=1pt,blue] (q)--(c);
\draw[->,line width=1pt,cyan] (c)--(r);
\end{tikzpicture}
\end{center}

随后仍要重复：
\[
c_0\longrightarrow c_1\longrightarrow c_2\longrightarrow\cdots\longrightarrow c_{L-1}.
\]

\smallbox{orange}{\parbox{0.92\linewidth}{
\textbf{轮询开销。}简单点积阶段需要 $Ld$ 次密文--明文逐槽乘法，并返回 $L$ 个“长度为 $m$ 的结果密文”。
当 $m=200,L=483$ 时，每个结果只使用 $200/4096=4.88\%$ 的槽位。}}

\newpage

% ========================================================================
\docsection{blue}{2. 查询向量的槽位扩展}

一个密文有 $S_{\rm slots}$ 个槽位。若一次处理 $m$ 个查询，则每个查询最多可以拥有
\[
T=\left\lfloor\frac{S_{\rm slots}}m\right\rfloor
\]
个连续位置。定义二维组合到一维槽位的映射：
\[
\boxed{\operatorname{slot}(q,t)=qT+t},
\qquad 0\le q<m,\quad 0\le t<T.
\]

对 $m=2,T=3$，一维槽位的语义不是六个无关数字，而是：
\[
\underbrace{(q_0,c_0),(q_0,c_1),(q_0,c_2)}_{q_0\text{ 的候选块}},
\quad
\underbrace{(q_1,c_0),(q_1,c_1),(q_1,c_2)}_{q_1\text{ 的候选块}}.
\]

为了让查询 $q_q$ 与 tile 中每个候选 $c_t$ 都相乘，查询的整行必须连续复制 $T$ 次：
\[
A^{\rm pack}=Q\otimes\vct 1_T\in\mathbb R^{mT\times d},
\qquad
A^{\rm pack}_{qT+t,f}=Q_{q,f}.
\]

具体例子中：
\[
Q=\begin{bmatrix}1&2\\3&4\end{bmatrix}
\quad\Longrightarrow\quad
A^{\rm pack}=\begin{bmatrix}
1&2\\1&2\\1&2\\
3&4\\3&4\\3&4
\end{bmatrix}_{6\times2}.
\]

逐特征读取 $A^{\rm pack}$ 的列，得到真正送入 CKKS 的一维向量：
\[
\begin{aligned}
\vct a_0&=[1,1,1,\;3,3,3]^{\mathsf T}\in\mathbb R^{mT},\\
\vct a_1&=[2,2,2,\;4,4,4]^{\mathsf T}\in\mathbb R^{mT}.
\end{aligned}
\]

\begin{center}
\slot{blue}{7mm}{1}\slot{blue}{7mm}{1}\slot{blue}{7mm}{1}
\slot{blue!65}{7mm}{3}\slot{blue!65}{7mm}{3}\slot{blue!65}{7mm}{3}
\hspace{4mm}$\leftarrow\vct a_0$\\[2mm]
\slot{blue}{7mm}{2}\slot{blue}{7mm}{2}\slot{blue}{7mm}{2}
\slot{blue!65}{7mm}{4}\slot{blue!65}{7mm}{4}\slot{blue!65}{7mm}{4}
\hspace{4mm}$\leftarrow\vct a_1$
\end{center}

\keybox{blue}{
\textbf{查询复制的必要性。}
若 $q_0$ 的特征值只出现一次，它在一次逐槽乘法中只能遇到一个候选值；
复制为 $T$ 份后，同一个查询值分别占据 $t=0,1,\ldots,T-1$ 的位置，才可能在一次 SIMD 运算中遇到 $T$ 个不同候选。

\textbf{查询扩展的计算含义。}
复制发生在加密前；随后仍然只按特征加密 $d$ 个向量
$\enc(\vct a_0),\ldots,\enc(\vct a_{d-1})$。
密文数量没有从 $d$ 变成 $dT$，只是单个密文的有效槽位从 $m$ 增加到 $mT$。
}

\newpage

% ========================================================================
\docsection{cyan}{3. 候选分块的槽位扩展与对齐}

B 从候选库中连续取 $T$ 行组成第 $j$ 个 tile：
\[
C^{(j)}=C[jT:(j+1)T,:]\in\mathbb R^{T\times d}.
\]
最后一个 tile 若只有 $v<T$ 个真实候选，就在末尾补 $T-v$ 行零；解包时只保留前 $v$ 列。

在例子中，$T=3$，候选 tile 为
\[
C^{(0)}=\begin{bmatrix}
10&20\\30&40\\50&60
\end{bmatrix}_{3\times2}.
\]

查询打包矩阵有 $mT$ 行，因此候选也必须变成 $mT$ 行。做法是把整个 tile 重复 $m$ 次：
\[
B^{(j)}_{\rm pack}=\vct 1_m\otimes C^{(j)}\in\mathbb R^{mT\times d},
\qquad
B^{(j)}_{qT+t,f}=C^{(j)}_{t,f}.
\]

具体展开为：
\[
B^{(0)}_{\rm pack}=\begin{bmatrix}
10&20\\30&40\\50&60\\
10&20\\30&40\\50&60
\end{bmatrix}_{6\times2}.
\]

它的逐特征明文向量为：
\[
\begin{aligned}
\vct b_0&=[10,30,50,\;10,30,50]^{\mathsf T}\in\mathbb R^{mT},\\
\vct b_1&=[20,40,60,\;20,40,60]^{\mathsf T}\in\mathbb R^{mT}.
\end{aligned}
\]

\keybox{cyan}{
\textbf{查询与候选的形状对应关系。}
$A^{\rm pack}$ 和 $B^{(j)}_{\rm pack}$ 都是 $mT\times d$。
在同一行 $r=qT+t$：
\[
A^{\rm pack}_{r,:}=q_q,
\qquad B^{(j)}_{{\rm pack},r,:}=c_{jT+t}.
\]
因此第 $r$ 行从一开始就代表同一对 $(q,t)$；不存在“查询在槽位 0、候选在槽位 100”的错位问题。
}

\begin{center}
\renewcommand{\arraystretch}{1.28}
\begin{tabular}{c c c c}
\toprule
槽位 $r$ & 组合语义 & $A^{\rm pack}_{r,:}$ & $B^{(0)}_{{\rm pack},r,:}$\\
\midrule
0 & $(q_0,c_0)$ & $[1,2]$ & $[10,20]$\\
1 & $(q_0,c_1)$ & $[1,2]$ & $[30,40]$\\
2 & $(q_0,c_2)$ & $[1,2]$ & $[50,60]$\\
3 & $(q_1,c_0)$ & $[3,4]$ & $[10,20]$\\
4 & $(q_1,c_1)$ & $[3,4]$ & $[30,40]$\\
5 & $(q_1,c_2)$ & $[3,4]$ & $[50,60]$\\
\bottomrule
\end{tabular}
\end{center}

\textcolor{muted}{注意：查询是“每一行重复 $T$ 次”；候选是“整个 $T\times d$ tile 重复 $m$ 次”。两种重复方向不同，但最终行索引完全一致。}

\newpage

% ========================================================================
\docsection{orange}{4. 逐槽乘加的正确性证明}

对每个特征 $f$，A 加密查询列 $\enc(\vct a_f)$，B 保留候选列明文 $\vct b_f$。
TenSEAL 执行逐槽乘法和逐槽加法：
\[
\vct z^{(j)}
=\sum_{f=0}^{d-1}\enc(\vct a_f)\odot\vct b_f
=\enc\!\left(\sum_{f=0}^{d-1}\vct a_f\odot\vct b_f\right)
\in\enc(\mathbb R^{mT}).
\]

\keybox{orange}{
\textbf{命题。}
对任意查询 $q$ 和 tile 内候选位置 $t$，令 $r=qT+t$，则
\[
z^{(j)}_r=\sum_{f=0}^{d-1}Q_{q,f}C^{(j)}_{t,f}
=\langle q_q,c_{jT+t}\rangle.
\]

\textbf{证明。}
由构造式可知 $a_{f,r}=Q_{q,f}$，同时 $b_{f,r}=C^{(j)}_{t,f}$。
逐槽运算不会混合不同下标，因此
\[
z^{(j)}_r
=\sum_f a_{f,r}b_{f,r}
=\sum_f Q_{q,f}C^{(j)}_{t,f}.
\]
这正是目标分数表中第 $(q,jT+t)$ 个元素。证毕。
}

把例子代入：
\[
\begin{aligned}
\vct z^{(0)}
&=\enc([1,1,1,3,3,3])\odot[10,30,50,10,30,50]\\
&\quad+\enc([2,2,2,4,4,4])\odot[20,40,60,20,40,60]\\
&=\enc([50,110,170,\;110,250,390]).
\end{aligned}
\]

按 query-major 规则恢复二维形状：
\[
\reshape_{m\times T}([50,110,170,110,250,390])
=\begin{bmatrix}
50&110&170\\
110&250&390
\end{bmatrix}
=QC^{\mathsf T}.
\]

项目最终还要减阈值并乘独立正随机掩码。对每个组合 $(q,t)$：
\[
\widetilde S_{q,t}
=R_{q,t}\left(\sum_fQ_{q,f}C_{t,f}-\tau\right),
\qquad R_{q,t}>0.
\]
这仍然是逐槽运算，不改变上面的槽位对应关系。

\smallbox{green}{\parbox{0.92\linewidth}{
\textbf{无旋转推论。}
每个点积都只在固定槽位 $r=qT+t$ 内沿特征维求和；不同槽位之间从不需要交换数据。
因此打包阶段和打分阶段都不需要 slot rotation。}}

\newpage

% ========================================================================
\docsection{cyan}{5. 项目两轮候选数据的构造}

上面的证明把 $C^{(j)}$ 写成一个普通候选 tile，这是槽位几何的核心。
在当前项目中，两轮的候选来源略有不同，但都服从同一个 $qT+t$ 布局。

\textbf{第一轮候选数据为聚类中心。}
B 有 $k=50$ 个聚类中心，每个中心是 $200$ 维特征：
\[
C_{\rm centroid}\in\mathbb R^{50\times200}.
\]
每次取 $T=20$ 个中心形成 $C^{(j)}\in\mathbb R^{20\times200}$，因此共有
\[
\left\lceil\frac{50}{20}\right\rceil=3\text{ 个 tiles}.
\]
这里同一组中心对所有查询相同，正好就是 $\vct1_m\otimes C^{(j)}$ 的直接情形。

\textbf{第二轮候选数据为所选聚类中的候选列。}
B 的库可以写成
\[
\mathcal B\in\mathbb R^{k\times L\times d},
\]
其中 $i$ 是聚类编号，$t$ 是该聚类内的候选位置，$f$ 是特征。
查询 $q$ 对应一个加密 one-hot selector $s_{q,i}$。因此它真正要比较的“有效候选”为
\[
\widehat C_{q,t,f}=\sum_{i=0}^{k-1}s_{q,i}\mathcal B_{i,t,f}.
\]

\keybox{cyan}{
\textbf{第二轮候选选择细节。}
第二轮并不是先把某个明文聚类公开选出来；B 将每个聚类候选的明文贡献与加密 selector 在同槽位相乘，
同态形成 $\widehat C_{q,t,f}$。因此不同查询可以在同一 tile 中选择不同聚类，但槽位仍然是 $qT+t$。
随后再与 $\enc(Q_{q,f})$ 做逐槽点积。
}

把第二轮单个 tile 的流程写成四步：

\begin{center}
\begin{tikzpicture}[node distance=5mm and 5mm,>=Stealth]
\node (db) [draw=cyan,rounded corners,fill=panel,minimum width=31mm,minimum height=15mm,align=center]
{B 候选库\\$\mathcal B[:,jT:(j+1)T,:]$};
\node (sel) [right=of db,draw=blue,rounded corners,fill=panel,minimum width=31mm,minimum height=15mm,align=center]
{加密 selector\\$\enc(s_{q,i})$};
\node (eff) [right=of sel,draw=orange,rounded corners,fill=panel,minimum width=35mm,minimum height=15mm,align=center]
{形成有效候选\\$\enc(\widehat C_{q,t,f})$};
\node (score) [right=of eff,draw=green,rounded corners,fill=panel,minimum width=34mm,minimum height=15mm,align=center]
{与查询逐槽点积\\$\widetilde S_{q,t}$};
\draw[->,line width=1pt,cyan] (db)--(sel);
\draw[->,line width=1pt,blue] (sel)--(eff);
\draw[->,line width=1pt,orange] (eff)--(score);
\end{tikzpicture}
\end{center}

第二轮 $L=483$，于是候选位置被切为
\[
[0,19],[20,39],\ldots,[460,479],[480,482]+17\text{ 行零填充},
\]
总计 $\lceil483/20\rceil=25$ 个 tiles。零填充槽位在解包时丢弃，不会被解释成真实候选。

\newpage

% ========================================================================
\docsection{orange}{6. 计算与通信开销比较}

设候选宽度为 $L$、特征维度为 $d$、查询数为 $m$、tile 宽度为 $T$。
先比较最核心的“普通候选点积”部分：

\begin{center}
\renewcommand{\arraystretch}{1.42}
\begin{tabularx}{0.98\linewidth}{>{\bfseries}p{38mm} X X}
\toprule
项目 & 修改前（逐候选轮询） & 修改后（二维 tile） \\
\midrule
一次结果的逻辑形状 & $m\times1$，只是一列 & $m\times T$，一次得到 $T$ 列 \\
查询输入 & $d$ 个长度 $m$ 的有效向量 & $d$ 个长度 $mT$ 的有效向量 \\
查询密文数量 & $d$ & $d$，\textcolor{green}{不增加} \\
B 端候选准备 & 每轮广播一个候选，$O(md)$ & 每个 tile 展开为 $mT\times d$，$O(mTd)$ \\
在线外层循环 & $L$ 次 & $\lceil L/T\rceil$ 次 \\
密文--明文乘法量 & $Ld$ & $\lceil L/T\rceil d$ \\
slot rotation & 不需要 & 不需要 \\
结果密文数 & $L$ & $\lceil L/T\rceil$ \\
\bottomrule
\end{tabularx}
\end{center}

\keybox{orange}{
\textbf{预处理开销。}
二维打包确实增加了加密前的槽位写入和 B 端每 tile 的明文展开：逻辑数据量从 $m$ 增加到 $mT$。
但 CKKS 多项式度和安全参数保持不变，查询仍是每个特征一个密文；B 端展开在本地完成，不增加 A 到 B 的候选通信。
换来的收益是：一次同态点积不再只产生 $m$ 个分数，而是产生 $mT$ 个分数。
}

当前参数代入：
\[
S_{\rm slots}=4096,\qquad m=200,\qquad
T=\left\lfloor\frac{4096}{200}\right\rfloor=20.
\]
\[
mT=200\times20=4000,\qquad
\text{槽位利用率}=\frac{4000}{4096}=97.66\%.
\]
\[
L=483:\qquad
483\text{ 次候选轮询}
\quad\Longrightarrow\quad
\left\lceil\frac{483}{20}\right\rceil=25\text{ 个 tiles}.
\]

\begin{center}
\begin{tikzpicture}
\node[draw=whiteish,fill=panel,rounded corners=3pt,inner sep=10pt,
      text=whiteish,text width=.90\linewidth,align=center]
{\normalsize\bfseries
\textcolor{whiteish}{结论}\hspace{3mm}
复制不是重复执行 $T$ 次同态计算；\\[1mm]
它是在加密前把 $m\times T$ 个 $(q,t)$ 组合排进同一密文，
让一次逐槽乘加自然产生 $mT$ 个分数。};
\end{tikzpicture}
\end{center}

因此实际加速应接近候选分块比例：
\[
\frac{483}{25}=19.32,
\qquad
\frac{2722.17\ \mathrm{s}}{144.42\ \mathrm{s}}\approx18.85.
\]
二者接近，说明收益主要来自每次并行处理约 $T=20$ 个候选，而不是修改 CKKS 算法、降低安全参数或省略原有打分步骤。

\end{document}
